\section{Discussion}

\subsection{Specification Drives Convergence}

The central finding of this study is not that domain knowledge constrains divergence --- it is that \textit{specification} does. Every design decision that is explicitly defined in the SCOPE.md brief converges perfectly across both runs: the scale (five feeders and one meta), the narrator voice (The Monk, with fixed stylistic rules), the domain (Age of Empires II), the audience (casual solo player), the format (printable Markdown), and the estimated solve time (2--3 hours) are identical across Hunt 1 and Hunt 2. Every design decision left to the pipeline's generative judgment --- round structure, meta mechanism family, answer word selection, clue formulations --- diverges substantially.

This is a stronger and more generalizable claim than the original hypothesis. The working hypothesis entering this study, drawn from design space theory \citep{simon1969sciences,lawson1997designers}, was that \textit{domain knowledge} would be the primary stabilizing force: because both hunts must draw on the same finite AoE2 knowledge base, the domain would constrain the output space and produce convergence. The results partially support this --- domain semantic gravity (discussed below) produces thematic coherence at the word-set level even where exact words diverge --- but the primary driver of convergence is not the domain. It is the specification. Decisions in SCOPE.md are stable; decisions not in SCOPE.md are variable, even when domain knowledge constrains the universe of valid answers to a relatively small set.

The practical implication for AI creative pipeline deployment is direct: to make a pipeline's outputs more predictable, specify more in the creative brief --- not by narrowing the creative space or constraining aesthetic choices, but by reducing the number of decisions the pipeline must resolve autonomously. The practitioners' question ``how much does the pipeline vary across runs?'' should be reframed as ``how much does my brief under-specify?'' In both observed runs, every under-specified decision was variable; every explicitly specified decision was stable in both runs. This reframing converts the same-input divergence problem from a property of the AI system (its stochasticity) into a property of the human-AI interface (the specification density of the creative brief). These recommendations should be treated as provisional guidance supported by two runs in a single domain, to be confirmed with N$\geq$5 replication before operational deployment.

Three candidate mechanisms could explain why specification produces convergence. First, a \emph{hard-filter} mechanism: specified decisions eliminate options that would violate the brief, leaving fewer valid choices. Second, a \emph{distributional-prior} mechanism: explicit specification shifts the model's token distribution toward brief-consistent options at each generation step. Third, an \emph{attention-anchor} mechanism: specified elements act as attention weights that steer downstream generation toward coherent continuations. These mechanisms predict different failure modes and should be distinguished in future work with targeted ablations.

The five comparison levels are not causally independent: Level~3 (answer word assignment) causally constrains Levels~4 and~5 by determining which content is available for mechanism selection and clue formulation. A hunt that assigned \textsc{grail} as its meta answer was necessarily constrained to produce a 5-letter anagram-compatible feeder set, which influenced mechanism choices at Level~4. This causal structure means that stabilising upstream decisions (Levels~1--3) has multiplicative stabilising effects downstream.

\subsection{Domain Semantic Gravity}

The Level 3 results reveal a phenomenon not anticipated by the original study design: \textit{domain semantic gravity}. The exact answer word sets diverge completely (Jaccard = 0.00 on string match; no word appears in both hunts). Yet 3 of the 5 feeder slots exhibit domain-adjacent word pairs across the two runs: (LOOM, FORGE) in the production-technology slot, (ONAGER, SIEGE) in the siege-warfare slot, (PATROL, MARCH) in the military-movement slot. The AoE2 domain does not force convergence to specific words --- it cannot, because the set of valid AoE2 vocabulary items is large and equally distributed across mechanism slots --- but it creates a semantic basin of attraction. Both runs fall into the same region of the AoE2 lexical space even when they choose different specific words within that region.

We term this domain semantic gravity: an observed pattern in this case study, consistent with an apparent tendency (N=2) of a domain-grounded creative generator to select outputs that cluster in a semantic subspace defined by the domain's vocabulary and conceptual structure, even across independent runs that do not share a single exact token; replication is required to establish it as a general property. Domain semantic gravity is a \textit{partial stabilizing force}: it does not produce stability at the string level, but it produces stability at the thematic level. This distinction is practically important for meta puzzle design. An experienced human constructor who needed exactly the word PATROL for a meta mechanism would not be helped by the knowledge that a different run would produce MARCH --- both fall in the military-verb region of AoE2 vocabulary, but they are different strings and cannot be substituted for one another in a mechanism that depends on specific letter positions or word length.

This observation also refines the practical recommendation from the previous subsection. Domain semantic gravity implies that multiple runs of the same under-specified pipeline will tend to produce answer word sets that are thematically coherent with the domain but lexically incompatible with one another. Practitioners who need specific answer word properties (e.g., a particular word length, a specific letter pattern, or a particular semantic category for meta purposes) cannot rely on domain semantic gravity to deliver them; they must either specify those properties explicitly in the brief, or run the pipeline multiple times and apply human selection from the resulting word sets.

\subsection{The Panel as a Stable Quality Signal}

Panel score stability is among this study's most practically important findings. If AI panel scores were highly sensitive to the arbitrary design choices made at Levels 3--5 --- answer words, extraction mechanisms, clue formulations --- then the panel's quality signal would be unreliable as a deployment gate: it would pass some runs and fail others based on the incidental luck of word selection rather than genuine differences in puzzle quality. The results show the opposite. The pass/fail verdict is stable across all scored puzzles in both hunts (all puzzles in both runs that received full panel review achieved a passing score). The normalized aggregate score is stable across the two runs (83\% and 86\% of maximum, a 3-point gap within expected inter-run noise). Clarity and Solvability --- the two dimensions most directly tied to the puzzle's structural correctness --- are nearly identical across runs (+0.2 and $-$0.3 on a 5-point scale, respectively).

In both runs, the panel's quality signal appears robust to the specific generative choices made at Levels 3--5, suggesting a quality dimension that may be invariant to those choices (N=2; replication required). The panel penalizes structurally broken extraction, ambiguous clue formulations, and domain-knowledge failures; it rewards clean aha moments, fair deductions, and thematic coherence. These properties attach to the puzzle's structure and the solver's cognitive path, not to the specific answer word or the particular technology chain chosen to instantiate that structure. Both runs produced structurally sound puzzles, and both runs were rewarded similarly by the panel.

The systematic upward shift in Elegance (+1.2) and Reading Reward (+1.2) from Hunt 1 to Hunt 2 is the one panel-level finding that is not merely stable noise. This shift is consistent with, but not conclusively attributable to, a genuine mechanism-level improvement; the double-weighting of these dimensions in Hunt 2's rubric and the change in reviewer panel are confounds that cannot be disentangled from this case study alone. Hunt 2's P1 (The Blacksmith's Chain) adopts a chain-following structure with deterministic extraction that passes the Computer Test cleanly and achieves near-perfect Riven Standard compliance; Hunt 1's structurally comparable puzzle (Economy Puzzle, V) required a redesign flag. This illustrates a key distinction: in both observed runs, the panel was stable enough to confirm that both hunts are deployable-quality, but sensitive enough to detect genuine quality differences when they exist.

\subsection{Limitations}

Three primary limitations bound the conclusions of this study. First, the empirical base is $N = 2$ runs. Two independent runs provide a case study of same-input divergence; they do not provide a statistical characterization of its distribution. The divergence patterns observed here --- Jaccard 0.00 at Level 3, 1/5 mechanism-type matches at Level 4, 3/5 domain-adjacent word pairs --- could be atypical of the pipeline's general behavior. A robust characterization of divergence structure requires at minimum $N \geq 10$ runs across the same domain, from which distribution statistics (mean, variance, and tail behavior of each divergence metric) can be estimated.

Second, the clean-room isolation between Hunt 1 and Hunt 2 was enforced procedurally --- the generating agent was instructed not to access Hunt 1 artifacts --- rather than architecturally. In a multi-agent system with separate process boundaries, architectural isolation would guarantee that Hunt 1's generated outputs could not influence Hunt 2's generation even if instruction following were imperfect. The procedural approach is sufficient for a controlled study but introduces a residual confound that architectural isolation would eliminate.

Third, only three of Hunt 2's five feeder puzzles were fully authored and panel-reviewed at the time of analysis. The panel score comparison covers three puzzle slots per hunt, which is adequate to observe the stability pattern at the direction level but insufficient to characterize the full score distribution or to identify any outlier puzzles in Hunt 2 that might fall below the passing threshold. The full comparison awaits completion of Hunt 2's P4 and P5.

Future work should replicate this study with $N = 10$ runs across at least three domains (AoE2, music, science fiction) to assess whether the specification-drives-convergence finding generalizes across domain types with different knowledge density and vocabulary structure. Automated comparison scripts (currently a manual coding step at Level 4) would reduce the cost of large-$N$ replication and enable near-real-time divergence monitoring in deployed creative pipelines.
